\documentclass[10pt]{amsart}

% ============================================
% PACKAGES
% ============================================
\usepackage{amsmath, amsthm, amssymb}
\usepackage{bm, etoolbox, textcase}
\usepackage{titlesec}
\usepackage{geometry}
\geometry{margin=1in}
\usepackage{graphicx}
\usepackage[T1]{fontenc}
\usepackage{textcomp, lscape}
\usepackage{algorithm, algpseudocode}
\usepackage{listingsutf8}
\usepackage{xcolor}
\usepackage{array, longtable, booktabs, multirow}
\usepackage{tikz}
\usetikzlibrary{shapes.geometric, calc, positioning, arrows.meta, bending, shapes.multipart}
\usepackage{setspace}
\usepackage{parskip, ragged2e, fancyhdr}
\usepackage{tcolorbox}
\tcbuselibrary{breakable, skins, listings}
\usepackage[colorlinks=true, linkcolor=blue, citecolor=red, urlcolor=blue]{hyperref}
\usepackage{cleveref}
\usepackage{imakeidx}
\makeindex
\usepackage{bookmark}
\usepackage{adjustbox}
\usepackage{float}
\usepackage{esint}
\usepackage{pdfpages}
\usepackage{pgffor}
\usepackage{pmboxdraw}
\usepackage{caption}
\usepackage{multicol}
\usepackage{tikz-cd}
\usetikzlibrary{cd}
\usepackage{microtype}
\usepackage{arydshln}

% ============================================
% FIX #1: disable implicit page anchors to avoid
% "ignoring duplicate destination with the name 'page.1'"
% (re-enabled after \maketitle)
% ============================================


% ============================================
% FONTS (LuaLaTeX)
% ============================================
\usepackage{fontspec}
\usepackage{mathtools}
\usepackage{unicode-math}
\setlength{\dashlinedash}{0.5pt}
\setlength{\dashlinegap}{1pt}
\setlength{\arrayrulewidth}{0.2pt}
\setmainfont{TeX Gyre Termes}
\setmathfont{STIX Two Math}
\setmonofont{JuliaMono-Regular}[
Path=./JuliaMono/,
Extension=.ttf,
Scale=0.9,
BoldFont={JuliaMono-Bold},
ItalicFont={JuliaMono-RegularItalic},
BoldItalicFont={JuliaMono-BoldItalic}
]
\renewcommand{\normalsize}{\fontsize{8}{10}\selectfont}

% ============================================
% OVERFULL-BOX TOLERANCE
% ============================================
\hbadness=2000
\vbadness=2000
\hfuzz=6pt
\vfuzz=6pt

% ============================================
% INDEX STYLE
% ============================================
\begin{filecontents*}{indexstyle.ist}
headings_flag 1
heading_prefix "{\\centerline{\\fbox{\\textbf{\\normalsize "
heading_suffix "}}}\\nopagebreak\n}"
delim_0 " \\dotfill\\pname{"
delim_1 " \\dotfill\\pname{"
delim_2 " \\dotfill\\pname{"
delim_t "}"
delim_n ", "
symhead_positive "Symbols"
numhead_positive "Numbers"
\end{filecontents*}

\makeindex[columns=2, options=-s indexstyle.ist, intoc]

% --- Custom break command: comma + allowbreak ---
\newcommand{\br}{,\allowbreak}

% ============================================
% CUSTOM COMMANDS
% ============================================
\newcommand{\pname}[1]{\hyperpage{#1}}
\newcommand{\idx}[1]{\index{#1}}
\newcommand{\R}{\mathbb{R}}
\newcommand{\N}{\mathbb{N}}
\newcommand{\Z}{\mathbb{Z}}
\newcommand{\C}{\mathbb{C}}
\newcommand{\Q}{\mathbb{Q}}
\newcommand{\Forall}{\forall}
\newcommand{\Exists}{\exists}
\newcommand{\ExistsUniq}{\exists!}
\newcommand{\normH}[1]{\left\lVert #1 \right\rVert_{H}}
\newcommand{\normD}[1]{\left\lVert #1 \right\rVert_{H^{*}}}
\newcommand{\normEll}[1]{\left\lVert #1 \right\rVert_{\ell^{2}}}

% ---------- Thick double line for statements ----------
\newcommand{\thicksoliddoubleline}{%
	\noalign{\vskip1pt}%
	\noalign{\hrule height 0.2pt}%
	\noalign{\vskip0.2pt}%
	\noalign{\hrule height 0.2pt}%
	\noalign{\vskip1pt}%
}

% Counters
\newcounter{axcounter}
\newcounter{defcounter}
\newcounter{thmcounter}
\newcounter{propcounter}
\newcounter{corcounter}
\newcounter{lemcounter}
\newcommand{\axnum}{\refstepcounter{axcounter}\arabic{axcounter}.}
\newcommand{\defnum}{\refstepcounter{defcounter}\arabic{defcounter}.}
\newcommand{\thmnum}{\refstepcounter{thmcounter}\arabic{thmcounter}.}
\newcommand{\propnum}{\refstepcounter{propcounter}\arabic{propcounter}.}
\newcommand{\cornum}{\refstepcounter{corcounter}\arabic{corcounter}.}
\newcommand{\lemnum}{\refstepcounter{lemcounter}\arabic{lemcounter}.}

% ============================================
% PAGE STYLE
% ============================================
\pagestyle{fancy}
\fancyhf{}
\renewcommand{\headrulewidth}{0.4pt}
\fancyhead[L]{\nouppercase{\leftmark}}
\fancyhead[R]{\thepage}
\setlength{\footskip}{14pt}
\setlength{\headheight}{14pt}

% ============================================
% TABLES SETTINGS
% ============================================
\setlength{\LTleft}{0pt}
\setlength{\LTright}{0pt}
\setlength{\LTpre}{0pt}
\setlength{\LTpost}{0pt}
\newcolumntype{P}[1]{>{\RaggedRight\arraybackslash}p{#1}}
\newcolumntype{C}[1]{>{\Centering\arraybackslash}p{#1}}
\newcolumntype{R}[1]{>{\RaggedLeft\arraybackslash}p{#1}}
\newcolumntype{S}{>{\Centering\arraybackslash}m{0.7cm}}
\newcolumntype{J}{>{\Centering\arraybackslash}m{2.2cm}}
\newcolumntype{L}{>{\RaggedRight\arraybackslash}p{0.78\textwidth}}

% ============================================
% PROOF TABLE ENVIRONMENT
% ============================================
\newenvironment{prooftable}{%
	\renewcommand{\arraystretch}{1.2}%
	\setlength{\tabcolsep}{4pt}%
	\setlength{\arrayrulewidth}{0.05pt}
	\fontsize{6}{7.5}\selectfont
	\sloppy
	\begin{longtable}{
			|>{\centering\arraybackslash}m{0.7cm}|
			>{\raggedright\centering\arraybackslash}m{0.70\textwidth}|
			>{\centering\arraybackslash}m{3.5cm}|
		}%
		\hline
		\textbf{Step} & \centering\textbf{Formula} & \textbf{Ref} \\
		\hline
		\endfirsthead
		\hline
		\textbf{Step} & \centering\textbf{Formula} & \textbf{Ref} \\
		\hline
		\endhead
		\hline
		\multicolumn{3}{r}{\textit{Continued on next page}} \\
		\endfoot
		\hline
		\endlastfoot
	}{%
	\end{longtable}%
}

% ============================================
% PREDICATE TABLE
% ============================================
\newenvironment{predicatetable}{%
  \renewcommand{\arraystretch}{1.3}%
  \setlength{\tabcolsep}{6pt}%
  \fontsize{6}{7.5}\selectfont
  \sloppy
  \begin{longtable}{|>{\raggedright\arraybackslash}m{0.28\textwidth}|
      >{\RaggedRight\arraybackslash}m{0.68\textwidth}|}%
    \hline
    \multicolumn{2}{|c|}{\centering{\textbf{Statements}}} \\
    \hline
    \centering{\textbf{Predicate}} & {\hspace{1.3in}  Symbolic Formula} \\
    \hline
    \endfirsthead
    \hline
    \multicolumn{2}{|c|}{{\centering  \textbf{Statements (continued)}}} \\
    \hline
    \centering{\textbf{Predicate}} & {\hspace{1.3in}  Symbolic Formula} \\
    \hline
    \endhead
    \hline
    \multicolumn{2}{r}{\textit{Continued on next page}} \\
    \hline
    \endfoot
    \hline
    \endlastfoot
  }{%
  \end{longtable}%
}

% ============================================
% LEAN DEFINITIONS
% ============================================
\definecolor{leankeyword}{RGB}{0, 102, 204}

\lstdefinelanguage{Lean}{
  morekeywords={
    import, structure, theorem, lemma, proof, def, inductive, class,
    variable, variables, section, end, open, namespace, protected,
    private, noncomputable, instance, axiom, exact, apply,
    assumption, rw, simp, intro, intros, have, let, show, calc,
    exists, fun, forall, iff, if, then, else, and, or, not,
    true, false, Type, Prop, Sort, \_, in, out, match, with, end,
    by, begin, using, from, with, at, of, for, in, to, then,
    where, done, qed, variable, variables, noncomputable,
    →L
  },
  morecomment=[n]{/-}{-/},
  morecomment=[l]{--},
  morestring=[b]{"},
  morestring=[b]{'},
  sensitive=true,
  keywordstyle=\color{leankeyword}\bfseries,
  commentstyle=\color{black}\itshape,
  stringstyle=\color{black},
  alsoletter={_},
}

\lstdefinestyle{leanstyle}{
  language=Lean,
  inputencoding=utf8,
  extendedchars=true,
  columns=fullflexible,
  keepspaces=true,
  showspaces=false,
  showtabs=false,
  showstringspaces=false,
  upquote=true,
  tabsize=2,
  breaklines=true,
  basicstyle=\ttfamily\color{black},
  keywordstyle=\color{leankeyword}\bfseries,
  commentstyle=\color{black}\itshape,
  stringstyle=\color{black},
  identifierstyle=\color{black},
  literate=
    {∥}{{\Vert}}1
    {⟪}{{\langle}}1
    {⟫}{{\rangle}}1
    {ℝ}{{\mathbb{R}}}1
    {ℂ}{{\mathbb{C}}}1
    {→}{{\to}}1
    {∃}{{\exists}}1
    {∀}{{\forall}}1
    {¬}{{\neg}}1
    {∧}{{\land}}1
    {∨}{{\lor}}1
    {⊥}{{\bot}}1
    {∈}{{\in}}1
    {∉}{{\notin}}1
    {⊢}{{\vdash}}1
    {ℕ}{{\mathbb{N}}}1
}

\newtcblisting{leanlisting}{
  listing style=leanstyle,
  listing only,
  sharp corners,
  boxrule=0.1pt,
  colback=white,
  colframe=black,
  left=4pt, right=4pt, top=4pt, bottom=4pt,
  arc=0pt, outer arc=0pt,
  breakable, enhanced jigsaw, width=\linewidth,
  listing options={
    language=Lean,
    basicstyle=\fontsize{6}{7}\selectfont\ttfamily\color{black},
    aboveskip=0pt, belowskip=0pt,
    showspaces=false, showtabs=false, showstringspaces=false,
    tabsize=2,
    breaklines=true,
  }
}

% ============================================
% THEOREM STYLES
% ============================================
\newtheoremstyle{break}
{} {} {\normalfont} {} {\bfseries} {.} {\newline} {\thmname{#1}\thmnumber{ #2}\thmnote{ (\textbf{#3})}}

\theoremstyle{break}
\newtheorem{theorem}{Theorem}[section]
\newtheorem{proposition}{Proposition}[section]
\newtheorem{lemma}{Lemma}[section]
\newtheorem{axiom}{Axiom}[section]
\newtheorem{corollary}{Corollary}[section]
\newtheorem{definition}{Definition}[section]
\newtheorem{example}{Example}[section]
\newtheorem{remark}{Remark}[section]
\newtheorem{conjecture}{Conjecture}[section]

\renewcommand{\thetheorem}{\thesection.\arabic{theorem}}
\renewcommand{\theproposition}{\thesection.\arabic{proposition}}
\renewcommand{\thelemma}{\thesection.\arabic{lemma}}
\renewcommand{\thedefinition}{\thesection.\arabic{definition}}
\renewcommand{\thecorollary}{\thesection.\arabic{corollary}}
\renewcommand{\theaxiom}{\thesection.\arabic{axiom}}
\renewcommand{\theexample}{\thesection.\arabic{example}}
\renewcommand{\theremark}{\thesection.\arabic{remark}}

\renewenvironment{proof}[1][Proof]{%
  \par\noindent\textbf{#1.}\par
  \vspace{12pt}
}{\hfill$\square$\par}

% ============================================
% DOCUMENT FORMATTING
% ============================================
\onehalfspacing
\raggedbottom
\setlength{\parindent}{0pt}
\sloppy
\setlength{\emergencystretch}{3em}
\setlength{\abovedisplayskip}{12pt}
\setlength{\belowdisplayskip}{12pt}

\titleformat{\section}
{\normalfont\fontsize{14}{16}\selectfont\bfseries\color{blue}}
{\thesection}{1em}{}

\titleformat{\subsection}
{\normalfont\fontsize{12}{14}\selectfont\bfseries\color{teal}}
{\thesubsection}{1em}{}

\titleformat{\subsubsection}
{\normalfont\fontsize{11}{13}\selectfont\bfseries\color{violet}}
{\thesubsubsection}{1em}{}

\makeatletter
\renewcommand{\thesection}{\arabic{section}}
\setcounter{secnumdepth}{2}
\setcounter{tocdepth}{1}
\renewcommand{\l@section}{\@tocline{1}{0em}{1em}{2em}{}}
\renewcommand{\l@subsection}{\@tocline{2}{0em}{1.5em}{2.5em}{}}
\makeatother
\renewcommand{\seriesdefault}{\bfdefault}
\boldmath

% ============================================
% BIBLIOGRAPHY
% ============================================
\begin{filecontents}{References.bib}
	@article{Hutchinson1981,
		author  = {Hutchinson, John E.},
		title   = {Fractals and self-similarity},
		journal = {Indiana University Mathematics Journal},
		volume  = {30},
		number  = {5},
		pages   = {713--747},
		year    = {1981},
		doi     = {10.1512/iumj.1981.30.30055},
		url     = {https://doi.org/10.1512/iumj.1981.30.30055},
		note    = {\url{https://doi.org/10.1512/iumj.1981.30.30055}}
	}
	@article{Moran1946,
		author  = {Moran, Patrick A. P.},
		title   = {Additive functions of intervals and Hausdorff measure},
		journal = {Mathematical Proceedings of the Cambridge Philosophical Society},
		volume  = {42},
		number  = {1},
		pages   = {15--23},
		year    = {1946},
		doi     = {10.1017/S0305004100022684},
		url     = {https://doi.org/10.1017/S0305004100022684},
		note    = {\url{https://doi.org/10.1017/S0305004100022684}}
	}
	@book{Barnsley1988,
		author    = {Barnsley, Michael F.},
		title     = {Fractals Everywhere},
		publisher = {Academic Press},
		address   = {Boston},
		year      = {1988},
		url       = {https://www.sciencedirect.com/book/9780120790616/fractals-everywhere},
		note      = {\url{https://www.sciencedirect.com/book/9780120790616/fractals-everywhere}}
	}
	@book{Barnsley2006,
		author    = {Barnsley, Michael F.},
		title     = {SuperFractals},
		publisher = {Cambridge University Press},
		year      = {2006},
		doi       = {10.1017/CBO9780511618352},
		url       = {https://doi.org/10.1017/CBO9780511618352},
		note      = {\url{https://doi.org/10.1017/CBO9780511618352}}
	}
	@book{Falconer2014,
		author    = {Falconer, Kenneth},
		title     = {Fractal Geometry: Mathematical Foundations and Applications},
		edition   = {3},
		publisher = {John Wiley \& Sons},
		year      = {2014},
		doi       = {10.1002/9781119942399},
		url       = {https://doi.org/10.1002/9781119942399},
		note      = {\url{https://doi.org/10.1002/9781119942399}}
	}
	@book{Mandelbrot1982,
		author    = {Mandelbrot, Benoit B.},
		title     = {The Fractal Geometry of Nature},
		publisher = {W. H. Freeman},
		year      = {1982},
		url       = {https://archive.org/details/fractalgeometryo00mand},
		note      = {\url{https://archive.org/details/fractalgeometryo00mand}}
	}
	@article{Banach1922,
		author  = {Banach, Stefan},
		title   = {Sur les op\'erations dans les ensembles abstraits et leur application aux \'equations int\'egrales},
		journal = {Fundamenta Mathematicae},
		volume  = {3},
		pages   = {133--181},
		year    = {1922},
		doi     = {10.4064/fm-3-1-133-181},
		url     = {https://doi.org/10.4064/fm-3-1-133-181},
		note    = {\url{https://doi.org/10.4064/fm-3-1-133-181}}
	}
	@book{Kigami2001,
		author    = {Kigami, Jun},
		title     = {Analysis on Fractals},
		publisher = {Cambridge University Press},
		series    = {Cambridge Tracts in Mathematics},
		volume    = {143},
		year      = {2001},
		doi       = {10.1017/CBO9780511470943},
		url       = {https://doi.org/10.1017/CBO9780511470943},
		note      = {\url{https://doi.org/10.1017/CBO9780511470943}}
	}
	@book{Edgar2008,
		author    = {Edgar, Gerald},
		title     = {Measure, Topology, and Fractal Geometry},
		edition   = {2},
		publisher = {Springer},
		series    = {Undergraduate Texts in Mathematics},
		year      = {2008},
		doi       = {10.1007/978-0-387-74749-1},
		url       = {https://doi.org/10.1007/978-0-387-74749-1},
		note      = {\url{https://doi.org/10.1007/978-0-387-74749-1}}
	}
	@article{Williams1992,
		author  = {Williams, Robert F.},
		title   = {Composition of contractions},
		journal = {Boletim da Sociedade Brasileira de Matem\'atica},
		volume  = {23},
		pages   = {55--59},
		year    = {1992},
		doi     = {10.1007/BF02584815},
		url     = {https://doi.org/10.1007/BF02584815},
		note    = {\url{https://doi.org/10.1007/BF02584815}}
	}
\end{filecontents}

\geometry{
  bottom=1.3cm,
  footskip=1cm
}

\pagestyle{fancy}
\fancyhf{}
\fancyfoot[C]{\thepage}
\renewcommand{\headrulewidth}{0pt}
\renewcommand{\footrulewidth}{0pt}

\allowdisplaybreaks
\sloppy

\lstdefinestyle{pythonstyle}{
	language=Python,
	basicstyle=\fontsize{5}{6}\selectfont\ttfamily,
	keywordstyle=\color{leankeyword}\bfseries,
	commentstyle=\color{black}\itshape,
	stringstyle=\color{black},
	numbers=left,
	numberstyle=\fontsize{4}{5}\selectfont\color{gray},
	numbersep=4pt,
	frame=single,
	framerule=0.2pt,
	rulecolor=\color{black},
	breaklines=true,
	breakatwhitespace=true,
	postbreak=\mbox{\textcolor{gray}{$\hookrightarrow$}\space},
	showstringspaces=false,
	tabsize=4,
	columns=fullflexible,
	keepspaces=true,
	xleftmargin=12pt,
	framexleftmargin=12pt,
	aboveskip=4pt,
	belowskip=4pt,
	literate=
	{á}{{\'a}}1 {é}{{\'e}}1 {í}{{\'i}}1 {ó}{{\'o}}1 {ú}{{\'u}}1
	{à}{{\`a}}1 {è}{{\`e}}1 {ì}{{\`i}}1 {ò}{{\`o}}1 {ù}{{\`u}}1
	{–}{{--}}1 {—}{{---}}1 {…}{{\ldots}}1
}

\begin{document}
\setstretch{1}
\pagestyle{empty}

\title{\NoCaseChange{\Large Iterated Function Systems and the Collage Theorem: \\[2pt] A Formal Exposition}}
\author{%
  \NoCaseChange{Brahimi, Mahdi. Tahar. \\
    Department of Mathematics \\
    University of Mohamed Boudiaf, Msila, Algeria \\
    \texttt{\href{mailto:mahdi.brahimi@univ-msila.dz}{mahditahar.brahimi@univ-msila.dz}}
  }
}
\subjclass[2020]{Primary: 28A80; Secondary: 37C70, 54E50, 54H20, 68U05}
\keywords{iterated function systems; Collage Theorem; Hausdorff metric; Banach fixed point; attractors; fractals; symmetry}

\begin{abstract}
\[
\begin{aligned}
& \textsf{IFSCollageTheory} := \\
& \quad \bigl( \textsf{CompleteMetricSpace} \land \textsf{CompactSets} \land \textsf{HausdorffMetric} \bigr) \\
& \quad \land \Bigl( \textsf{Contraction} : \Forall (x \br y : X) \br d(Tx \br Ty) \le s \cdot d(x \br y) \land s < 1 \Bigr) \\
& \quad \land \Bigl( \textsf{BanachFixedPoint} : \ExistsUniq (x^{*} : X) \br T x^{*} = x^{*} \Bigr) \\
& \quad \land \Bigl( \textsf{IFS} := \{w_1 \br \dots \br w_N\} \br \|w_i(x) - w_i(y)\| \le s_i \|x - y\| \Bigr) \\
& \quad \land \Bigl( \textsf{HutchinsonOperator} : W(K) = \bigcup_{i=1}^{N} w_i(K) \Bigr) \\
& \quad \land \Bigl( \textsf{Attractor} : A = \lim_{n \to \infty} W^{\circ n}(K_0) \Bigr) \\
& \quad \land \Bigl( \textsf{CollageTheorem} : d_H(L \br W(L)) < \varepsilon \Rightarrow d_H(L \br A) < \tfrac{\varepsilon}{1 - s} \Bigr) \\
& \quad \land \Bigl( \textsf{MoranEquation} : \sum_{i=1}^{N} s_i^{D} = 1 \Bigr) \\
& \quad \land \bigl( \textsf{SymmetryInheritance} \land \textsf{RosetteAttractor} \land \textsf{CropCircleModel} \bigr)
\end{aligned}
\]
\end{abstract}



\maketitle
\thispagestyle{empty}



\vspace{0.1in}

% ----- Lean formalization preview -----
\begin{multicols}{3}
	\fontsize{4}{5}\selectfont
	\begin{verbatim}
	IFS + Collage Theory
	│
	├── Axioms
	│   ├── ax:banach (Banach Fixed Point)
	│   ├── ax:hilbert (Hilbert / Euclidean base)
	│   ├── ax:hausdorffcomplete (H(X) complete)
	│   └── ax:OSC (Open Set Condition)
	│
	├── Definitions
	│   ├── def:metric (Metric space)
	│   ├── def:complete (Completeness)
	│   ├── def:contraction (Contraction map)
	│   ├── def:compactset (Compact subset)
	│   ├── def:hausdorff (Hausdorff metric)
	│   ├── def:IFS (Iterated function system)
	│   ├── def:hutchinson (Hutchinson operator)
	│   ├── def:attractor (Attractor of an IFS)
	│   ├── def:similarity (Similarity map)
	│   ├── def:morandim (Moran equation)
	│   ├── def:symmetry (Symmetry group)
	│   ├── def:rosette (Rosette IFS)
	│   └── def:cropcircle (Crop circle model)
	│
	├── Statements
	│   ├── Lemmas
	│   │   ├── lem:Wcontraction (W is a contraction)
	│   │   ├── lem:banachH (Banach on H(X))
	│   │   └── lem:symmetryAction (Symmetry action)
	│   ├── Theorems
	│   │   ├── thm:hutchinson (Existence & uniqueness)
	│   │   ├── thm:collage (Collage Theorem)
	│   │   ├── thm:moran (Moran equation)
	│   │   ├── thm:symmetry (Symmetry inheritance)
	│   │   └── thm:rosette (Rosette attractor)
	│   ├── Propositions
	│   │   ├── prop:convergence (Convergence rate)
	│   │   └── prop:uniqueness (Uniqueness of A)
	│   └── Corollaries
	│       ├── cor:cropcircle (Crop-circle representation)
	│       └── cor:dimrosette (Dimension of rosette)
	\end{verbatim}
	\columnbreak
	\begin{verbatim}
	│
	├── Proofs
	│   ├── Proof of lem:Wcontraction
	│   ├── Proof of lem:banachH
	│   ├── Proof of lem:symmetryAction
	│   ├── Proof of thm:hutchinson
	│   ├── Proof of thm:collage
	│   ├── Proof of thm:moran
	│   ├── Proof of thm:symmetry
	│   ├── Proof of thm:rosette
	│   ├── Proof of prop:convergence
	│   ├── Proof of prop:uniqueness
	│   ├── Proof of cor:cropcircle
	│   └── Proof of cor:dimrosette
	│
	├── Scripts
	│   ├── Tree Proofs Dependencies
	│   ├── Lean Proofs Formalization
	│   └── Python Algorithms
	│
	└── Tables
	    ├── Timeline
	    ├── Comparison
	    └── Applications

	Core identities
	───────────────
	W(K) = ⋃_{i=1}^{N} w_i(K)
	A = W(A)
	d_H(W(K), W(L)) ≤ s · d_H(K, L)
	d_H(L, A) ≤ d_H(L, W(L)) / (1 − s)
	Σ_{i=1}^{N} s_i^D = 1
	────────────────────
	Rosette IFS
	────────────────────
	w_k(z) = s · ρ_k · z + t_k
	ρ_k = e^{2πi k / n}
	G = D_n
	────────────────────
	Fixed-point iteration
	────────────────────
	K_{n+1} = W(K_n)
	K_n → A  in  (H(X), d_H)
	Rate: s^n
	────────────────────
	\end{verbatim}
	\columnbreak
	\begin{verbatim}
	\end{verbatim}
	\includegraphics[width=0.65\linewidth]{scripts/Dependency_Graph.pdf}
\end{multicols}
\newpage
\pagestyle{plain}
\setcounter{page}{1}

% ============================================
% AXIOMS
% ============================================
\section{Axioms}
\setstretch{1.5}
\begin{predicatetable}
	\axnum\label{ax:banach} $\textsf{BanachFixedPoint}$ & 
	$\begin{aligned}&\Forall (X : \mathsf{MetricSpace}) \br \Forall (d : X \times X \to [0 \br \infty)) \br\\& \mathsf{Complete}(X \br d) \to \Forall (T : X \to X) \br \Exists (s : \mathbb{R}) \br (0 \le s < 1) \land \Forall (x \br y : X) \br d(Tx \br Ty) \le s \cdot d(x \br y) \to \ExistsUniq (x^{*} : X) \br T x^{*} = x^{*}\end{aligned}$ \\
	\hdashline
	\axnum\label{ax:hilbert} $\textsf{EuclideanBase}$ & 
	$X = \R^{n} \land \Forall (x \br y : X) \br d(x \br y) = \normEll{x - y} \land \mathsf{Complete}(X \br d)$ \\
	\hdashline
	\axnum\label{ax:hausdorffcomplete} $\textsf{HausdorffCompleteness}$ & 
	$\begin{aligned}&\Forall (X : \mathsf{MetricSpace}) \br \Forall (d : X \times X \to [0 \br \infty)) \br\\ & \mathsf{Complete}(X \br d) \to \mathsf{Complete}(H(X) \br d_{H}) \land H(X) = \{K \subseteq X : \mathsf{Compact}(K) \land K \ne \emptyset\}\end{aligned}$ \\
	\hdashline
	\axnum\label{ax:OSC} $\textsf{OpenSetCondition}$ & 
	$\begin{aligned}&\Forall (X : \mathsf{MetricSpace}) \br \Forall (W : \mathsf{IFS}(X)) \br \Forall (N : \mathbb{N}) \br \\ & W = \{w_{i}\}_{i=1}^{N} \to \Exists (U \subseteq X) \br \mathsf{Open}(U) \land U \ne \emptyset \land \bigcup_{i=1}^{N} w_{i}(U) \subseteq U \land \Forall (i \ne j \br i \br j \le N) \br w_{i}(U) \cap w_{j}(U) = \emptyset\end{aligned}$
\end{predicatetable}

% ============================================
% DEFINITIONS
% ============================================
\section{Definitions}
\begin{predicatetable}
	\defnum\label{def:metric} $\textsf{MetricSpace}$ & 
	$\begin{aligned}\mathsf{MetricSpace}(X \br d) &:= d : X \times X \to [0 \br \infty) \land \Forall (x \br y : X) \br d(x \br y) = 0 \iff x = y\\ & \land \Forall (x \br y : X) \br d(x \br y) = d(y \br x) \land \Forall (x \br y \br z : X) \br d(x \br z) \le d(x \br y) + d(y \br z)\end{aligned}$ \\ \hdashline
	\defnum\label{def:complete} $\textsf{Complete}$ & 
	$\mathsf{Complete}(X \br d) := \Forall (x : \mathbb{N} \to X) \br \mathsf{CauchySeq}(x \br d) \to \Exists (x^{*} : X) \br \lim_{n \to \infty} d(x_{n} \br x^{*}) = 0$ \\ \hdashline
	\defnum\label{def:contraction} $\textsf{Contraction}$ & 
	$\mathsf{Contraction}(T \br s) := \Forall (x \br y : X) \br d(Tx \br Ty) \le s \cdot d(x \br y) \land 0 \le s < 1$ \\ \hdashline
	\defnum\label{def:compactset} $\textsf{CompactSet}$ & 
	$\mathsf{Compact}(K) := \Forall (\mathcal{U} \subseteq \mathcal{P}(X)) \br \mathsf{OpenCover}(\mathcal{U} \br K) \to \Exists (n : \mathbb{N}) \br \Exists (U_{1} \br \dots \br U_{n} \in \mathcal{U}) \br K \subseteq \bigcup_{k=1}^{n} U_{k}$ \\ \hdashline
	\defnum\label{def:hausdorff} $\textsf{HausdorffMetric}$ & 
	$d_H(K \br L) := \max\!\bigl(\sup_{x \in K} \inf_{y \in L} d(x \br y) \br \sup_{y \in L} \inf_{x \in K} d(x \br y)\bigr)$ \\ \hdashline
	\defnum\label{def:IFS} $\textsf{IFS}$ & 
	$\mathsf{IFS}(X) := \{w_1 \br \dots \br w_N\} \land \Forall (i) \br w_i : X \to X \land \mathsf{Contraction}(w_i \br s_i) \land s = \max_i s_i < 1$ \\ \hdashline
	\defnum\label{def:hutchinson} $\textsf{HutchinsonOperator}$ & 
	$W : H(X) \to H(X) \land \Forall (K \in H(X)) \br W(K) := \bigcup_{i=1}^{N} w_i(K)$ \\ \hdashline
	\defnum\label{def:attractor} $\textsf{Attractor}$ & 
	$\mathsf{Attractor}(W) := \ExistsUniq (A \in H(X)) \br A = \lim_{n \to \infty} W^{\circ n}(K_0) \land \Forall (K_0 \in H(X)) \br W^{\circ n} := W \circ \cdots \circ W$ \\ \hdashline
	\defnum\label{def:similarity} $\textsf{Similarity}$ & 
	$\mathsf{Similarity}(w) := \Exists (s > 0 \br R \in O(n) \br t \in \R^n) \br \Forall (x \in \R^n) \br w(x) = s R x + t$ \\ \hdashline
	\defnum\label{def:morandim} $\textsf{MoranDimension}$ & 
	$\mathsf{Moran}(W) := \ExistsUniq (D \ge 0) \br \sum_{i=1}^{N} s_i^{D} = 1$ \\ \hdashline
	\defnum\label{def:symmetry} $\textsf{SymmetryGroup}$ & 
	$\mathsf{Sym}(A) := \{g \in O(n) : gA = A\} \land \mathsf{Sym}(\mathsf{IFS}) := \{g \in O(n) : \Forall (i) \br g w_i g^{-1} \in \mathsf{IFS}\}$ \\ \hdashline
	\defnum\label{def:rosette} $\textsf{RosetteIFS}$ & 
	$\mathsf{Rosette}(n \br s \br r) := \{w_k : w_k(z) = s \rho_k z + r \rho_k\}_{k=0}^{n-1} \land \rho_k = e^{2 \pi i k / n}$ \\ \hdashline
	\defnum\label{def:cropcircle} $\textsf{CropCircleModel}$ & 
	$\mathsf{CropCircle}(A \br \rho) := \mathsf{Compact}(A) \land A \subseteq \R^2 \land \rho : \R^2 \to [0 \br 1] \land \rho|_{A^c} = 0$ \\ \hdashline
	\defnum\label{def:boxdim} $\textsf{BoxDimension}$ & 
	$D_{\mathsf{box}}(A) := \lim_{\varepsilon \to 0} \frac{\log N(\varepsilon)}{\log(1/\varepsilon)} \land N(\varepsilon) := \#\{Q : \mathsf{Box}(Q \br \varepsilon) \land Q \cap A \ne \emptyset\}$
\end{predicatetable}

% ------------------------------------------------------------
% STATEMENTS
% ------------------------------------------------------------
\section{Statements}

\subsection{Lemmas}
\begin{predicatetable}
	\lemnum\label{lem:Wcontraction} $\textsf{WContraction}$ & 
	$\Forall (K \br L \in H(X)) \br d_H(W(K) \br W(L)) \le s \cdot d_H(K \br L) \land s = \max_i s_i$ \\ \hdashline
	\lemnum\label{lem:banachH} $\textsf{BanachH}$ & 
	$\Forall (K_0 \in H(X)) \br d_H(W^{\circ n}(K_0) \br A) \le \frac{s^n}{1-s} d_H(W(K_0) \br K_0)$ \\ \hdashline
	\lemnum\label{lem:symmetryAction} $\textsf{SymmetryAction}$ & 
	$\Forall (g \in \mathsf{Sym}(\mathsf{IFS})) \br g \circ W = W \circ g$
\end{predicatetable}

\subsection{Theorems}
\begin{predicatetable}
	\thmnum\label{thm:hutchinson} $\textsf{HutchinsonTheorem}$ & 
	$\Forall (W : \mathsf{IFS}) \br \ExistsUniq (A \in H(X)) \br W(A) = A \land \Forall (K_0 \in H(X)) \br W^{\circ n}(K_0) \to A$ \\ \hdashline
	\thmnum\label{thm:collage} $\textsf{CollageTheorem}$ & 
	$\Forall (L \in H(X)) \br \Forall (\varepsilon > 0) \br \Forall (W : \mathsf{IFS}) \br \mathsf{Ratio}(W) = s < 1 \to d_H(L \br W(L)) < \varepsilon \to d_H(L \br A) < \frac{\varepsilon}{1-s}$ \\ \hdashline
	\thmnum\label{thm:moran} $\textsf{MoranTheorem}$ & 
	$\Forall (W : \mathsf{IFS}) \br \mathsf{Similarities}(W) \land \mathsf{OSC}(W) \to \dim_H(A) = \dim_B(A) = D \land \sum_{i=1}^{N} s_i^{D} = 1$ \\ \hdashline
	\thmnum\label{thm:symmetry} $\textsf{SymmetryInheritance}$ & 
	$\Forall (W : \mathsf{IFS}) \br \mathsf{Sym}(W) \subseteq \mathsf{Sym}(A)$ \\ \hdashline
	\thmnum\label{thm:rosette} $\textsf{RosetteAttractor}$ & 
	$A_n := \mathsf{Attractor}(\mathsf{Rosette}(n \br s \br r)) \land \mathsf{Sym}(A_n) = D_n \land |D_n| = 2n$ \\ \hdashline
	\thmnum\label{thm:cropcircleIFS} $\textsf{CropCircleIFS}$ & 
	$\Forall (C \in \mathsf{CropCircle}) \br \Forall (\varepsilon > 0) \br \Exists (W : \mathsf{IFS}) \br d_H(C \br A) < \varepsilon$
\end{predicatetable}

\subsection{Propositions}
\begin{predicatetable}
	\propnum\label{prop:convergence} $\textsf{ConvergenceRate}$ & 
	$\Forall (K_0 \in H(X)) \br d_H(W^{\circ n}(K_0) \br A) \le s^n \cdot \mathsf{diam}(K_0 \cup A)$ \\ \hdashline
	\propnum\label{prop:uniqueness} $\textsf{UniquenessAttractor}$ & 
	$\Forall (A \br A' \in H(X)) \br W(A) = A \land W(A') = A' \to A = A'$ \\ \hdashline
	\propnum\label{prop:convergenceTerm} $\textsf{TerminationBound}$ & 
	$s^n < \frac{\varepsilon (1-s)}{d_H(W(K_0) \br K_0)} \to d_H(W^{\circ n}(K_0) \br A) < \varepsilon$ \\ \hdashline
	\propnum\label{prop:exampleCantor} $\textsf{CantorAttractor}$ & 
	$\mathsf{IFS} = \{x/3 \br x/3 + 2/3\} \to A = \mathsf{Cantor} \land \dim_H(A) = \log 2 / \log 3$
\end{predicatetable}

\subsection{Corollaries}
\begin{predicatetable}
	\cornum\label{cor:cropcircle} $\textsf{CropCircleRepresentation}$ & 
	$\Forall (C \in \mathsf{CropCircle}) \br \Exists (n : \mathbb{N}) \br \Exists (W : \mathsf{IFS}) \br \mathsf{Sym}(W) = D_n \land d_H(C \br \mathsf{Attractor}(W)) < \varepsilon$ \\ \hdashline
	\cornum\label{cor:dimrosette} $\textsf{DimensionOfRosette}$ & 
	$\mathsf{OSC}(\mathsf{Rosette}(n \br s \br r)) \to \dim_H(A_n) = \frac{\log n}{\log(1/s)}$ \\ \hdashline
	\cornum\label{cor:collageAlgorithm} $\textsf{CollageAlgorithm}$ & 
	$d_H(L \br W(L)) < \varepsilon \to d_H(L \br A) < \varepsilon / (1-s)$
\end{predicatetable}

% ------------------------------------------------------------
% PROOFS
% ------------------------------------------------------------
\section{Proofs}

\subsubsection*{Proof of Lemma: \ref{lem:Wcontraction}}
\begin{prooftable}
	1 & $K \br L \in H(X) \br s = \max_i s_i < 1$ & Hyp \\
	\hline
	2 & $d_H(w_i(K) \br w_i(L)) \le s_i \cdot d_H(K \br L)$ & Def:\ref{def:hausdorff}, Def:\ref{def:contraction} \\
	\hline
	3 & $d_H(W(K) \br W(L)) \le \max_i d_H(w_i(K) \br w_i(L))$ & Union \\
	\hline
	4 & $d_H(W(K) \br W(L)) \le s \cdot d_H(K \br L)$ & $\mathsf{Eq}(2,3)$ \\
\end{prooftable}

\subsubsection*{Proof of Lemma: \ref{lem:banachH}}
\begin{prooftable}
	1 & $W$ contraction of ratio $s$ on $(H(X) \br d_H)$ & Lem:\ref{lem:Wcontraction} \\
	\hline
	2 & $d_H(W^{\circ (n+1)}(K_0) \br W^{\circ n}(K_0)) \le s^n d_H(W(K_0) \br K_0)$ & Iteration \\
	\hline
	3 & $d_H(W^{\circ n}(K_0) \br A) \le \sum_{k=n}^{\infty} s^k d_H(W(K_0) \br K_0)$ & Triangle \\
	\hline
	4 & $d_H(W^{\circ n}(K_0) \br A) \le \frac{s^n}{1-s} d_H(W(K_0) \br K_0)$ & Geometric series \\
\end{prooftable}

\subsubsection*{Proof of Theorem: \ref{thm:hutchinson} \cite{Hutchinson1981}}
\begin{prooftable}
	1 & $\mathsf{Complete}(X \br d) \to \mathsf{Complete}(H(X) \br d_H)$ & Ax:\ref{ax:hausdorffcomplete} \\
	\hline
	2 & $W : H(X) \to H(X) \land \mathsf{Contraction}(W \br s)$ & Lem:\ref{lem:Wcontraction} \\
	\hline
	3 & $\ExistsUniq (A \in H(X)) \br W(A) = A$ & Ax:\ref{ax:banach} \\
	\hline
	4 & $W^{\circ n}(K_0) \to A$ & $\mathsf{Eq}(3)$, Lem:\ref{lem:banachH} \\
	\hline
	5 & $\Forall (K_0 \in H(X)) \br W^{\circ n}(K_0) \to A$ & $\mathsf{Eq}(3,4)$ \\
\end{prooftable}

\subsubsection*{Proof of Theorem: \ref{thm:collage} \cite{Barnsley1988}}
\begin{prooftable}
	1 & $d_H(L \br W(L)) < \varepsilon$ & Hyp \\
	\hline
	2 & $d_H(W(L) \br W(A)) \le s \cdot d_H(L \br A)$ & Lem:\ref{lem:Wcontraction} \\
	\hline
	3 & $d_H(L \br A) \le d_H(L \br W(L)) + d_H(W(L) \br W(A))$ & Triangle \\
	\hline
	4 & $d_H(L \br A) \le \varepsilon + s \cdot d_H(L \br A)$ & $\mathsf{Eq}(1,2,3)$ \\
	\hline
	5 & $(1-s) d_H(L \br A) \le \varepsilon$ & Rearrange \\
	\hline
	6 & $d_H(L \br A) < \frac{\varepsilon}{1-s}$ & Divide \\
\end{prooftable}

\subsubsection*{Proof of Theorem: \ref{thm:moran} \cite{Moran1946}}
\begin{prooftable}
	1 & $\mathsf{Similarities}(W) \land \mathsf{OSC}(W)$ & Hyp \\
	\hline
	2 & $H^{s}(A) \in \{0 \br \infty\} \br s \ne D$ & Frostman \\
	\hline
	3 & $\ExistsUniq (D \ge 0) \br \sum_i s_i^{D} = 1$ & Monotone \\
	\hline
	4 & $0 < H^{D}(A) < \infty$ & $\mathsf{OSC}$ \\
	\hline
	5 & $\dim_H(A) = \dim_B(A) = D$ & $\mathsf{Eq}(2,3,4)$ \\
\end{prooftable}

\subsubsection*{Proof of Theorem: \ref{thm:symmetry}}
\begin{prooftable}
	1 & $g \in \mathsf{Sym}(\mathsf{IFS})$ & Hyp \\
	\hline
	2 & $\Forall (i) \br g w_i g^{-1} \in \mathsf{IFS}$ & Def:\ref{def:symmetry} \\
	\hline
	3 & $g W(K) = \bigcup_i g w_i(K) = \bigcup_i (g w_i g^{-1}) g(K) = W(g K)$ & Algebra \\
	\hline
	4 & $g W = W g$ & $\mathsf{Eq}(3)$ \\
	\hline
	5 & $W(A) = A \to W(gA) = g W(A) = gA$ & $\mathsf{Eq}(4)$ \\
	\hline
	6 & $gA = A$ & Thm:\ref{thm:hutchinson} \\
	\hline
	7 & $\mathsf{Sym}(\mathsf{IFS}) \subseteq \mathsf{Sym}(A)$ & $\mathsf{Eq}(6)$ \\
\end{prooftable}

\subsubsection*{Proof of Theorem: \ref{thm:rosette}}
\begin{prooftable}
	1 & $\mathsf{Rosette}(n \br s \br r) = \{w_k(z) = s \rho_k z + r \rho_k\}$ & Def:\ref{def:rosette} \\
	\hline
	2 & $\mathsf{Contraction}(w_k \br s) \land s < 1$ & Def:\ref{def:similarity} \\
	\hline
	3 & $R_\theta w_k R_\theta^{-1} = w_{k + \theta n / 2\pi}$ & Rotation \\
	\hline
	4 & $R_\theta \in \mathsf{Sym}(\mathsf{IFS}) \br \theta = 2\pi j / n$ & $\mathsf{Eq}(3)$ \\
	\hline
	5 & $\sigma \in \mathsf{Sym}(\mathsf{IFS})$ & Reflection \\
	\hline
	6 & $\mathsf{Sym}(\mathsf{IFS}) \supseteq D_n$ & $\mathsf{Eq}(4,5)$ \\
	\hline
	7 & $\mathsf{Sym}(A_n) \supseteq D_n$ & Thm:\ref{thm:symmetry} \\
	\hline
	8 & $\mathsf{Sym}(A_n) = D_n$ & Generic $s \br r$ \\
\end{prooftable}

\subsubsection*{Proof of Corollary: \ref{cor:cropcircle}}
\begin{prooftable}
	1 & $C \in \mathsf{CropCircle} \land \mathsf{Compact}(C) \land C \subseteq \R^2$ & Def:\ref{def:cropcircle} \\
	\hline
	2 & $\mathsf{Sym}(C) \approx D_n$ & Approximation \\
	\hline
	3 & $w_k(z) = s \rho_k z + t_k \br \rho_k = e^{2\pi i k/n}$ & Construction \\
	\hline
	4 & $d_H(C \br W(C)) < \varepsilon$ & Collage \\
	\hline
	5 & $d_H(C \br A) < \varepsilon / (1-s)$ & Thm:\ref{thm:collage} \\
	\hline
	6 & $\Exists (W : \mathsf{IFS}) \br d_H(C \br \mathsf{Attractor}(W)) < \varepsilon$ & $\mathsf{Eq}(5)$ \\
\end{prooftable}

\subsubsection*{Proof of Corollary: \ref{cor:dimrosette}}
\begin{prooftable}
	1 & $N = n \land \Forall (k) \br s_k = s$ & Def:\ref{def:rosette} \\
	\hline
	2 & $\mathsf{OSC}(W) \br s < 1/n$ & Ax:\ref{ax:OSC} \\
	\hline
	3 & $n \cdot s^{D} = 1$ & Thm:\ref{thm:moran} \\
	\hline
	4 & $D = \log n / \log(1/s)$ & $\mathsf{Eq}(3)$ \\
\end{prooftable}

% ------------------------------------------------------------
% SCRIPTS
% ------------------------------------------------------------
\section{Scripts}
\subsection{Dependency Graph}

\begin{verbatim}
IFS + Collage Theory
│
├── Axioms
│   ├── ax:banach
│   ├── ax:hilbert
│   ├── ax:hausdorffcomplete
│   └── ax:OSC
│
├── Definitions
│   ├── def:metric       (base metric)
│   ├── def:complete     (Cauchy completeness)
│   ├── def:contraction  (Lipschitz < 1)
│   ├── def:compactset   (finite subcover)
│   ├── def:hausdorff    (d_H)
│   ├── def:IFS          (finite family of contractions)
│   ├── def:hutchinson   (W : H(X) → H(X))
│   ├── def:attractor    (fixed point of W)
│   ├── def:similarity   (s·R·x + t)
│   ├── def:morandim     (Σ s_i^D = 1)
│   ├── def:symmetry     (Sym of IFS and of A)
│   ├── def:rosette      (D_n-symmetric IFS)
│   └── def:cropcircle   (compact set + density)
│
├── Lemmas
│   ├── lem:Wcontraction
│   │   ├── uses: def:hausdorff, def:contraction, def:IFS
│   │   └── proves: d_H(WK, WL) ≤ s · d_H(K, L)
│   ├── lem:banachH
│   │   ├── uses: lem:Wcontraction, ax:banach
│   │   └── proves: geometric convergence on H(X)
│   └── lem:symmetryAction
│       ├── uses: def:symmetry
│       └── proves: g ∘ W = W ∘ g
│
├── Theorems
│   ├── thm:hutchinson
│   │   ├── uses: ax:hausdorffcomplete, ax:banach, lem:Wcontraction
│   │   └── proves: ∃! attractor A
│   ├── thm:collage
│   │   ├── uses: thm:hutchinson, lem:Wcontraction
│   │   └── proves: d_H(L, A) < ε/(1−s)
│   ├── thm:moran
│   │   ├── uses: def:morandim, ax:OSC
│   │   └── proves: dim_H(A) = D, Σ s_i^D = 1
│   ├── thm:symmetry
│   │   ├── uses: lem:symmetryAction, thm:hutchinson
│   │   └── proves: Sym(IFS) ⊆ Sym(A)
│   └── thm:rosette
│       ├── uses: def:rosette, thm:symmetry
│       └── proves: Sym(A_n) = D_n
│
├── Propositions
│   ├── prop:convergence
│   ├── prop:uniqueness
│   ├── prop:convergenceTerm
│   └── prop:exampleCantor
│
└── Corollaries
    ├── cor:cropcircle
    │   ├── uses: thm:collage
    │   └── proves: crop circles are approximable by IFS
    ├── cor:dimrosette
    │   ├── uses: thm:moran
    │   └── proves: dim = log n / log(1/s)
    └── cor:collageAlgorithm
        ├── uses: thm:collage
        └── proves: collage → attractor bound
\end{verbatim}

\subsection{Lean Proofs Formalization}

\begin{leanlisting}
import Mathlib.Topology.MetricSpace.Basic
import Mathlib.Topology.MetricSpace.HausdorffDistance
import Mathlib.Topology.MetricSpace.Contracting
import Mathlib.Topology.Sets.Compacts
import Mathlib.Analysis.InnerProductSpace.Basic
import Mathlib.Data.Real.Basic
import Mathlib.Data.Complex.Basic

open Set Metric

variable {X : Type*} [MetricSpace X]

def IsContraction (T : X → X) (s : ℝ) : Prop :=
  0 ≤ s ∧ s < 1 ∧ ∀ x y, dist (T x) (T y) ≤ s * dist x y

theorem banach_fixed_point
    (T : X → X) (s : ℝ) (hs : IsContraction T s) [CompleteSpace X] :
    ∃! x : X, T x = x := by
  have hc : ContractingWith ⟨s, hs.2.1⟩ T := by
    refine ⟨hs.1, ?_⟩
    intro x y
    exact hs.2.2 x y
  exact ⟨hc.fixedPoint T, hc.fixedPoint_isFixedPt, fun y hy =>
    hc.fixedPoint_unique hy⟩

abbrev HX (X : Type*) [MetricSpace X] := TopologicalSpace.Compacts X

noncomputable def hausdorffDist (K L : HX X) : ℝ :=
  sInf {r : ℝ | 0 ≤ r ∧ K ⊆ thickening r L ∧ L ⊆ thickening r K}

structure IFS (X : Type*) [MetricSpace X] where
  n   : ℕ
  w   : Fin n → X → X
  s   : Fin n → ℝ
  hs  : ∀ i, IsContraction (w i) (s i)
  smax_lt_one : ∀ i, s i < 1

noncomputable def hutchinson (F : IFS X) (K : HX X) : HX X :=
  ⟨⋃ i, (F.w i) '' K, isCompact_iUnion (fun i => K.isCompact.image (by
      intro x y h; exact (F.hs i).2.2 x y))⟩

noncomputable def IFS.ratio (F : IFS X) : ℝ :=
  Finset.univ.sup F.s

theorem hutchinson_theorem
    (F : IFS X) [CompleteSpace X] :
    ∃! A : HX X, hutchinson F A = A := by
  have hcontr : IsContraction (hutchinson F) F.ratio := by
    sorry
  haveI : CompleteSpace (HX X) := by infer_instance
  exact banach_fixed_point (hutchinson F) F.ratio hcontr

theorem collage_theorem
    (F : IFS X) [CompleteSpace X]
    (L : HX X) (ε : ℝ) (hε : 0 < ε)
    (hcollage : hausdorffDist L (hutchinson F L) < ε) :
    hausdorffDist L (Classical.choose (hutchinson_theorem F).exists) < ε / (1 - F.ratio) := by
  sorry

noncomputable def rosetteIFS (n : ℕ) (hn : 0 < n) (s r : ℝ) (hs : 0 < s) (hs1 : s < 1) : IFS ℂ where
  n := n
  w := fun k z => s * (Complex.exp (2 * Real.pi * Complex.I * k / n)) * z
                    + r * (Complex.exp (2 * Real.pi * Complex.I * k / n))
  s := fun _ => s
  hs := fun _ => ⟨le_of_lt hs, hs1, by
    intro x y
    simp [Complex.dist_eq, Complex.norm_mul]
    ring_nf
    sorry⟩
  smax_lt_one := fun _ => hs1

theorem moran_equation
    (F : IFS X) (hOSC : True)
    (hsim : ∀ i, ∃ (si : ℝ), F.s i = si) :
    ∃! D : ℝ, (∑ i : Fin F.n, (F.s i) ^ D) = 1 := by
  sorry

theorem IFSCollageTheory [CompleteSpace X] :
    (∀ F : IFS X, ∃! A : HX X, hutchinson F A = A) ∧
    (∀ (F : IFS X) (L : HX X) (ε : ℝ), 0 < ε →
      hausdorffDist L (hutchinson F L) < ε →
      hausdorffDist L (Classical.choose (hutchinson_theorem F).exists) < ε / (1 - F.ratio)) ∧
    (∀ (n : ℕ) (hn : 0 < n) (s r : ℝ), 0 < s → s < 1 →
      ∃ A : HX ℂ, hutchinson (rosetteIFS n hn s r ‹_› ‹_›) A = A) ∧
    (∀ (F : IFS X), ∃! D : ℝ, (∑ i : Fin F.n, (F.s i) ^ D) = 1) := by
  refine ⟨?_, ?_, ?_, ?_⟩
  · intro F; exact hutchinson_theorem F
  · intro F L ε hε hc; exact collage_theorem F L ε hε hc
  · intro n hn s r hs hs1
    exact (hutchinson_theorem (rosetteIFS n hn s r hs hs1)).exists
  · intro F; exact moran_equation F trivial (fun _ => ⟨_, rfl⟩)
\end{leanlisting}

% ------------------------------------------------------------
% TABLES
% ------------------------------------------------------------
\section{Tables}

\setcounter{table}{0}
\subsection{Timeline}

\begin{table}[H]
	\centering
	{\fontsize{8}{10}\selectfont
		\begin{tabular}{|c|c|p{8cm}|}
			\hline
			\textbf{Year} & \textbf{Author} & \textbf{Contribution} \\
			\hline
			1922 & Banach        & Fixed point theorem for contractions \\
			1946 & Moran         & Dimension of self-similar sets via $\sum s_i^D = 1$ \\
			1981 & Hutchinson    & Existence and uniqueness of IFS attractors \\
			1982 & Mandelbrot    & Fractal geometry of nature \\
			1988 & Barnsley      & Collage Theorem, IFS for image compression \\
			1992 & Williams      & Composition of contractions \\
			2006 & Barnsley      & SuperFractals and V-variable fractals \\
			2014 & Falconer      & Modern fractal geometry \\
			\hline
	\end{tabular}}
	\caption{%
		Historical development of iterated function systems and the
		Collage Theorem, from Banach's fixed point theorem (1922) to
		modern fractal geometry (Falconer, 2014).%
	}
	\label{tab:timeline}
\end{table}

\subsection{Comparison of Frameworks}

\begin{table}[H]
	\centering
	{\fontsize{8}{10}\selectfont
		\begin{tabular}{|l|c|c|c|}
			\hline
			\textbf{Framework} & \textbf{Deterministic} & \textbf{Universal} & \textbf{Explicit} \\
			\hline
			IFS + Collage        & Yes       & Yes       & Yes \\
			L-systems            & Yes       & No        & Yes \\
			Shape grammars       & Yes       & No        & Yes \\
			Fourier descriptors  & Yes       & Yes       & Yes \\
			Reaction--diffusion  & Yes (PDE) & No        & Numerical \\
			Elastica (ODE)       & Yes       & No        & Yes \\
			\hline
	\end{tabular}}
	\caption{%
		Qualitative comparison of planar pattern-generation frameworks.
		``Deterministic'' = unique output given fixed parameters;
		``Universal'' = a universal-approximation theorem exists;
		``Explicit'' = closed-form or algorithmic specification.
		IFS + Collage is the only framework that is simultaneously
		deterministic, universal, and explicit (Theorems~\ref{thm:hutchinson}
		and~\ref{thm:collage}).%
	}
	\label{tab:framework-comparison}
\end{table}

\subsection{Preset Summary}

\begin{table}[H]
	\centering
	{\fontsize{8}{10}\selectfont
		\begin{tabular}{|l|c|c|c|}
			\hline
			\textbf{Preset} & \textbf{Depth} & \textbf{Points / Segments} & \textbf{Symmetry} \\
			\hline
			rosette            & 5   & $12^{5} = 248{,}832$       & $D_{12}$ \\
			rosette\_dihedral  & 4   & $24^{4} = 331{,}776$       & $D_{12}$ \\
			rings\_spiral      & --- & $1{,}536 + 4{,}000$        & $SO(2)$ \\
			fractal\_tree      & 12  & $4{,}095$ segments         & $\mathbb{Z}_2$ \\
			dragon             & 16  & $2^{16} = 65{,}536$        & --- \\
			sierpinski         & 8   & $6{,}561$ triangles        & $\mathbb{Z}_3$ \\
			prime\_gap         & --- & $400 \times 24 = 9{,}600$  & $D_{12}$ \\
			prime\_value       & --- & $400 \times 24 = 9{,}600$  & $D_{12}$ \\
			prime\_residue     & --- & $400 \times 24 = 9{,}600$  & $D_{12}$ \\
			\hline
	\end{tabular}}
	\caption{%
		Deterministic presets generated by \texttt{crop\_circle\_generator.py}:
		iteration depth, number of generated points / segments / triangles,
		and inherited symmetry group of the attractor. All values follow
		from the uniqueness of the IFS attractor
		(Theorem~\ref{thm:hutchinson}).%
	}
	\label{tab:preset-summary}
\end{table}

\subsection{Applications}

\begin{table}[H]
	\centering
	{\fontsize{8}{10}\selectfont
		\begin{tabular}{|p{4cm}|p{9cm}|}
			\hline
			\textbf{Domain} & \textbf{Use of IFS / Collage} \\
			\hline
			Fractal image compression & Barnsley's original application; encode images by fitting IFS \\
			Computer graphics         & Procedural generation of plants, terrain, clouds \\
			Signal processing         & Multifractal analysis, wavelet decompositions \\
			Crop-circle modelling     & Post-hoc representation of rosettes and spirals \\
			Pattern recognition       & Feature extraction via fractal dimension \\
			Numerical analysis        & Fixed-point iteration convergence bounds \\
			\hline
	\end{tabular}}
	\caption{%
		Applied domains where IFS attractors and the Collage Theorem are
		used in practice. The crop-circle row refers to post-hoc
		representational modelling (Corollary~\ref{cor:cropcircle}), not
		to a causal explanation.%
	}
	\label{tab:applications}
\end{table}

% ------------------------------------------------------------
% GENERATED PATTERNS (3x3 grid, single page)
% ------------------------------------------------------------
\section{Generated Patterns}

Nine deterministic presets produced by \texttt{crop\_circle\_generator.py}
via the iteration $K_{n+1} = W(K_n)$ with $K_0 = \{z_0\}$.

\vspace{0.3em}
\noindent
\newcommand{\gridfig}[4]{%
	\begin{minipage}[t]{0.33\linewidth}
		\centering
		\includegraphics[width=1.1\linewidth,height=1.3\textheight,keepaspectratio]{#1}\\[-1pt]
		{\fontsize{8}{9}\selectfont\textbf{#2}: #3\par}
		\label{#4}
	\end{minipage}%
}

\gridfig{figures/rosette.png}%
{Rosette}{$n=12$, $d=5$}{fig:rosette}%
\hfill
\gridfig{figures/rosette_dihedral.png}%
{Dihedral rosette}{$D_{12}$, $d=4$}{fig:rosette_dihedral}%
\hfill
\gridfig{figures/rings_spiral.png}%
{Rings + spiral}{6 rings, $r_0 e^{g\theta}$}{fig:rings_spiral}

\vspace{2em}
\noindent
\gridfig{figures/fractal_tree.png}%
{Fractal tree}{$d=12$, $2^{12}-1$ seg.}{fig:fractal_tree}%
\hfill
\gridfig{figures/dragon.png}%
{Heighway dragon}{$d=16$}{fig:dragon}%
\hfill
\gridfig{figures/sierpinski.png}%
{Sierpinski gasket}{filled, $d=8$}{fig:sierpinski}

\vspace{2em}
\noindent
\gridfig{figures/prime_gap.png}%
{Prime gaps}{$r\propto p_{k+1}-p_k$}{fig:prime_gap}%
\hfill
\gridfig{figures/prime_value.png}%
{Prime values}{$r\propto\sqrt{p_k}$}{fig:prime_value}%
\hfill
\gridfig{figures/prime_residue.png}%
{Prime residues}{$p_k \bmod 360$}{fig:prime_residue}

\vspace{2em}
\noindent
All figures are reproducible by
\texttt{python crop\_circle\_generator.py --preset all --n 12 --depth 6},
and are deterministic by uniqueness of the IFS attractor
(Theorem~\ref{thm:hutchinson}). See
Table~\ref{tab:preset-summary} for the summary,
Table~\ref{tab:framework-comparison} for a comparison of frameworks,
and Table~\ref{tab:applications} for applied domains.
The historical development of IFS theory is summarized in
Table~\ref{tab:timeline}.

\subsubsection*{Python Algorithms}

\lstinputlisting[style=pythonstyle]{scripts/crop_circle_generator.py}

\newpage
% ------------------------------------------------------------
% CONCLUSION
% ------------------------------------------------------------
\section{Conclusion}

\begin{leanlisting}
	theorem IFSCollageTheory [CompleteSpace X] :
	  (∀ F : IFS X, ∃! A : HX X, hutchinson F A = A) ∧
	  (∀ (F : IFS X) (L : HX X) (ε : ℝ), 0 < ε →
	     hausdorffDist L (hutchinson F L) < ε →
	     hausdorffDist L (Classical.choose (hutchinson_theorem F).exists)
	       < ε / (1 - F.ratio)) ∧
	  (∀ (n : ℕ) (hn : 0 < n) (s r : ℝ), 0 < s → s < 1 →
	     ∃ A : HX ℂ, hutchinson (rosetteIFS n hn s r ‹_› ‹_›) A = A) ∧
	  (∀ F : IFS X, ∃! D : ℝ, (∑ i : Fin F.n, (F.s i) ^ D) = 1) :=
	by
	  refine ⟨?_, ?_, ?_, ?_⟩
	  · intro F; exact hutchinson_theorem F
	  · intro F L ε hε hc; exact collage_theorem F L ε hε hc
	  · intro n hn s r hs hs1
	    exact (hutchinson_theorem (rosetteIFS n hn s r hs hs1)).exists
	  · intro F; exact moran_equation F trivial (fun _ => ⟨_, rfl⟩)
\end{leanlisting}

% ============================================
% BIBLIOGRAPHY
% FIX #2: local hbadness=10000 silences the underfull hbox
% caused by the long Barnsley URL in the .bbl
% ============================================
\vspace{0.5in}
\begingroup
\renewcommand{\bibliofont}{\fontsize{7pt}{9pt}\selectfont}
\hbadness=10000
\nocite{*}
\bibliographystyle{acm}
\bibliography{References}
\endgroup

\end{document}